\section{Étude de faisabilité}

\subsection{But}
Pour réaliser une transmission vidéo robuste face aux obstacles, notre stratégie repose sur le forçage 
d'un MCS bas et l'utilisation de la technique de correction d'erreurs LDPC\index{LDPC}. 
L'objectif de cette étude de faisabilité est de déterminer si ces paramètres, généralement gérés 
dynamiquement par le matériel, sont configurables manuellement sur les modules Wi-Fi HaLow Morse 
Micro MM8108\index{Wi-Fi HaLow!MM8108} utilisés pour ce projet.

\subsection{Méthodologie d'investigation}
Dans un premier temps, l'analyse de la documentation constructeur et des interfaces utilisateur 
(\textit{user-friendly}) d'OpenWrt\index{OpenWrt} a révélé que si certains paramètres basiques 
sont accessibles (comme le choix du canal ou de la largeur de bande), l'activation forcée du LDPC 
ou le verrouillage strict du MCS ne sont pas exposés par les utilitaires standards 
(tels que \texttt{iw}, \texttt{morse\_cli} ou \texttt{hostapd}). L'investigation s'est donc orientée 
vers le reverse engineering et l'analyse du code source des pilotes fournis par Morse Micro, tant pour 
l'environnement FreeRTOS (Émetteur embarqué) que pour l'environnement Linux (Routeur relais).

\subsection{Faisabilité côté Émetteur (FreeRTOS / STM32)}
L'émetteur (TX), chargé de récupérer la vidéo et d'émettre les ondes radio depuis le drone, repose sur 
un microcontrôleur STM32\index{STM32} exploitant le système d'exploitation temps réel 
FreeRTOS\index{FreeRTOS}. L'investigation du kit de développement logiciel (\texttt{mm-iot-sdk}) 
a permis d'identifier les leviers d'action suivants :
\begin{itemize}
    \item \textbf{Contrôle du MCS :} L'analyse des bibliothèques a mis en évidence la fonction interne 
    \texttt{mmwlan\_ate\_override\_rate\_control()} issue du fichier d'en-tête \texttt{mmwlan.h}. 
    Cette fonction de test (ATE) permet de verrouiller dynamiquement la bande passante et le MCS 
    directement dans le code source de l'application en C.
    \item \textbf{Activation du LDPC :} L'inspection des configurations internes, notamment celles 
    liées au service \texttt{wpa\_supplicant} (\texttt{config\_ssid.h}), indique que la fonctionnalité 
    LDPC est activée par défaut par le pilote lors de la négociation de la connexion.
\end{itemize}

\subsection{Faisabilité côté Récepteur (Linux / OpenWrt)}
Le relais radio, configuré en point d'accès (Access Point, AP) au sol pour réceptionner le flux, 
repose sur une architecture Linux (OpenWrt). L'analyse du code source du pilote noyau 
(\textit{kernel driver}) a révélé les éléments suivants :
\begin{itemize}
    \item \textbf{Contrôle du MCS :} Les fichiers sources \texttt{rc.c} et \texttt{command.c} contiennent 
    des structures permettant de désactiver l'algorithme d'adaptation de débit \newline  
    (\texttt{enable\_fixed\_rate}) et d'imposer un MCS fixe (par exemple, MCS 2). Des tests de 
    compilation croisée (\textit{cross-compilation}) du noyau ont confirmé que ces modifications 
    logicielles peuvent être compilées sous forme de module (\texttt{morse.ko}) et injectées à 
    chaud sur le routeur cible.
    \item \textbf{La limite du forçage LDPC :} Initialement, une modification directe des en-têtes de 
    paquets MAC a été explorée. L'étude du fichier \texttt{skbq.c} laissait supposer la possibilité 
    d'intercepter les trames pour manipuler le registre \texttt{morse\_ratecode} 
    (en forçant les bits B17 et B18 du champ SIG-1). Cependant, les expérimentations ont démontré qu'il 
    est techniquement impossible d'appliquer et de maintenir ces paramètres matériels bas niveau tant 
    que le routeur n'est pas formellement connecté à un client. 
    
    La preuve formelle d'une injection manuelle du LDPC n'a pas pu être clairement établie
    et a donc été écartée de l'implémentation finale. 
    Néanmoins, la présence active des paramètres liés au LDPC et au MCS 10 dans le code source du pilote 
    confirme la capacité de la puce à utiliser cette correction d'erreur en interne. Nous partons donc 
    du postulat techniquement fondé que le contrôleur radio l'exploite de manière native.
\end{itemize}

\subsection{Conclusion de la faisabilité}
L'étude démontre que l'objectif de manipulation de la couche physique est réalisable sur les points 
critiques. S'il s'est avéré impraticable de forger manuellement chaque en-tête MAC pour imposer le 
LDPC en raison de la machine d'états du standard Wi-Fi, l'accès au code source du SDK FreeRTOS et 
des pilotes Linux permet d'appliquer les surcharges (\textit{overrides}) nécessaires pour verrouiller 
l'utilisation d'un MCS bas. Ces découvertes valident la faisabilité d'une liaison radio déterministe, 
prérequis absolu pour notre système vidéo.